\section{Trabalhos Relacionados}

Pesquisas sobre elasticidade de recursos tratam camadas e problemas distintos. \citeonline{pasin2013infaestrutura} estudam o ajuste autonômico de réplicas de servidores em uma infraestrutura distribuída sujeita a variações de carga. \citeonline{maciel2020gerenciamento} propõe uma arquitetura de gerenciamento e orquestração de fatias elásticas que integra recursos de nuvem e conectividade SDN. Lima Filho~(\citeyear{limafilho2026controle}) avalia o compartilhamento de blocos de recursos em redes de acesso 5G Open RAN por meio de uma xApp. Esses trabalhos tratam, respectivamente, da elasticidade de servidores, da orquestração de fatias e dos recursos de rádio; nenhum deles trata a seleção de caminhos em topologias de transporte em malha considerada aqui.

No fatiamento elástico de redes, \citeonline{gomes2020strategies} analisam as estratégias PETIC, PSR, SE e CSS sob variações temporais da demanda. Este trabalho toma o PETIC como referência exata para o custo aditivo implementado e acrescenta ACO e Q-Learning como buscas aproximadas desse mesmo custo, além do SA com função objetivo que inclui equilíbrio de utilização e extensão de caminho. A contribuição é experimental: disponibiliza um simulador comum e reprodutível e estima, com cenários independentes e comparações pareadas, como essas implementações se comportam sob as mesmas requisições e restrições de capacidade. Não se propõe um novo algoritmo de alocação nem uma avaliação geral das meta-heurísticas ou do aprendizado por reforço.
